\documentclass[11pt,a4paper]{article}

\usepackage[utf8]{inputenc}
\usepackage[T1]{fontenc}
\usepackage[margin=2.5cm]{geometry}
\usepackage{booktabs}
\usepackage{array}
\usepackage{xcolor}
\usepackage{hyperref}
\usepackage{amsmath}
\usepackage{amssymb}
\usepackage{enumitem}
\usepackage{multirow}
\usepackage{caption}
\usepackage{parskip}

\definecolor{darkblue}{RGB}{0,50,120}
\hypersetup{colorlinks=true,linkcolor=darkblue,urlcolor=darkblue}

\newcolumntype{L}[1]{>{\raggedright\arraybackslash}p{#1}}
\newcolumntype{C}[1]{>{\centering\arraybackslash}p{#1}}
\newcolumntype{R}[1]{>{\raggedleft\arraybackslash}p{#1}}

\title{%
  \textbf{Betweenness Centrality Prediction with Hyper-Connected GNNs}\\[0.4em]
  \large Project Progress Report
}
\author{Research Log Synthesis}
\date{March 11, 2026}

\begin{document}
\maketitle
\tableofcontents
\newpage

%------------------------------------------------------------
\section{Research Question}
%------------------------------------------------------------

Can deep GNNs predict betweenness centrality (BC) inductively,
and do HC / mHC / mHC-lite improve depth scaling and
out-of-distribution (OOD) ranking quality
compared with strong deep-GNN baselines?

\textbf{Primary evaluation metrics:}
\begin{itemize}[noitemsep]
  \item Ranking: Kendall $\tau$, Spearman $\rho$
  \item Top-K: Precision@K, NDCG@K
\end{itemize}

\textbf{Protocol (inductive):} train on small synthetic graphs
(100 nodes, ER/BA/SBM), evaluate on larger unseen graphs
(500 nodes, same families = ID; WS/RGG families = OOD).

%------------------------------------------------------------
\section{Stage 1 --- Full 5-Seed Depth Sweep (2026-03-06)}
%------------------------------------------------------------

\subsection{Setup}
\begin{table}[h]
\centering
\begin{tabular}{L{4.5cm} L{7.5cm}}
\toprule
Seeds & 5 (0--4) \\
Models & 19 variants: GCN, SAGE, GAT, GIN, GCNII, APPNP, JKNet,
         HC, mHC, mHC-lite \\
Depths & $L \in \{2,4,8,16,32\}$ \\
Train graphs & 160 (100 nodes, ER/BA/SBM) \\
Val / Test ID / OOD & 30 / 30 / 30 (500 nodes) \\
Budget & 50 epochs, patience 10, LR $10^{-3}$, WD $10^{-5}$ \\
\bottomrule
\end{tabular}
\end{table}

\subsection{Best points from the frozen sweep}
\begin{table}[h]
\centering
\begin{tabular}{L{4.5cm} L{2.5cm} C{1.2cm} R{3.5cm}}
\toprule
Category & Model & $L$ & OOD Kendall \\
\midrule
Best overall OOD     & \texttt{hc\_gcnii}  & 2 & $0.3824 \pm 0.0253$ \\
Best mHC point       & \texttt{mhc\_gcn}   & 2 & $0.3812 \pm 0.0238$ \\
Strongest GCN gain   & \texttt{hc\_gcn}    & 4 & $0.3809 \pm 0.0386$ \\
Best deep baseline   & \texttt{appnp}      & 32 & $0.3544 \pm 0.0258$ \\
\bottomrule
\end{tabular}
\end{table}

\subsection{Conclusions}
HC / mHC improve OOD consistently; gains concentrate at $L \leq 8$.
No ultra-deep scaling claim is supported.
Two families retained: \textbf{GCN} (clearest mHC signal)
and \textbf{GCNII} (deep-backbone stress test).

%------------------------------------------------------------
\section{Stage 2 --- Targeted Hyperparameter Search (2026-03-06)}
%------------------------------------------------------------

Optuna proxy search (15 trials, validation-only objective):

\begin{itemize}
  \item \textbf{GCN Trial 3} (obj. 0.8605): LR $2.26\times10^{-4}$,
        WD $2.10\times10^{-5}$, dropout 0.1, $n_\text{str}=2$,
        $\tau=0.2$, iters 20.
  \item \textbf{GCNII Trial 11} (obj. 0.8564): LR $1.07\times10^{-3}$,
        WD $4.34\times10^{-5}$, $n_\text{str}=4$, $\tau=0.1$,
        GCNII $\alpha=0.1$, $\theta=0.5$.
\end{itemize}

GCNII failure at L16 persists after tuning.
GCN + mHC is the strongest candidate.

%------------------------------------------------------------
\section{Stage 3 --- Multi-Seed Confirmation (2026-03-07)}
%------------------------------------------------------------

90 runs per family (3 candidate trials $\times$ 3 depths $\times$ 5 seeds).

\textbf{GCN:} mHC-GCN consistently gains on OOD Kendall at L4/L8
(strongest: $+0.071$ at t3/L4). Single negative case: t13/L2 ($-0.006$).

\textbf{GCNII:} mHC helps at L2/L8 but degrades at L16
(deltas: $-0.039$, $-0.012$, $-0.047$).

\textit{Decision: GCN family only from this point.}

%------------------------------------------------------------
\section{Stage 4 --- GCN Ablation Study (2026-03-07)}
%------------------------------------------------------------

\subsection{Results at L8 (tuned GCN regime, 5 seeds)}
\begin{table}[h]
\centering
\begin{tabular}{L{3.5cm} L{3.5cm} R{2.5cm} R{2.5cm}}
\toprule
Block & Variant & Val Kendall & OOD Kendall \\
\midrule
Structure & gcn\_ref          & 0.8332 & 0.3574 \\
Structure & mhc\_gcn\_ref     & 0.8909 & 0.3888 \\
Structure & \textbf{hc\_gcn\_ref} & \textbf{0.9174} & \textbf{0.4054} \\
Structure & dynamic only      & 0.8882 & 0.3904 \\
Structure & static only       & 0.8832 & 0.3855 \\
Structure & $n=4$ streams     & 0.8809 & 0.3791 \\
\midrule
Sinkhorn  & $\tau=0.05$       & 0.8873 & 0.3923 \\
Sinkhorn  & $\tau=0.10$       & 0.8932 & 0.3973 \\
Sinkhorn  & $\tau=0.20$       & 0.8952 & 0.3911 \\
Sinkhorn  & iters $=10$       & 0.8886 & 0.3906 \\
\bottomrule
\end{tabular}
\end{table}

\subsection{Interpretation}
HC-GCN dominates. No mHC variant exceeds HC regardless of $\tau$ tuning.
$n=2$ streams is optimal. Dynamic and static components contribute equally.
The useful mechanism is flexible multi-stream routing, \emph{not} the
doubly-stochastic Sinkhorn constraint.

%------------------------------------------------------------
\section{Stage 5 --- Final GCN Family Comparison (2026-03-09)}
%------------------------------------------------------------

\begin{table}[h]
\centering
\caption{Final GCN family --- test OOD results (5 seeds)}
\begin{tabular}{L{3cm} C{0.8cm} R{2.2cm} R{2.2cm} R{2.2cm} R{2.2cm}}
\toprule
Variant & $L$ & Kendall & Spearman & P@top5\% & NDCG@top5\% \\
\midrule
gcn            & 4 & $0.302\pm0.045$ & $0.428\pm0.064$ & 0.302 & 0.530 \\
gcn            & 8 & $0.360\pm0.021$ & $0.506\pm0.028$ & 0.315 & 0.580 \\
\midrule
hc\_gcn        & 4 & $0.393\pm0.031$ & $0.543\pm0.041$ & 0.341 & 0.589 \\
hc\_gcn        & 8 & $\mathbf{0.406\pm0.041}$ & $\mathbf{0.558\pm0.054}$ & 0.351 & 0.588 \\
\midrule
mhc\_gcn       & 4 & $0.374\pm0.043$ & $0.521\pm0.055$ & 0.329 & 0.581 \\
mhc\_gcn       & 8 & $0.396\pm0.038$ & $0.548\pm0.049$ & 0.343 & $\mathbf{0.590}$ \\
\midrule
mhc\_lite\_gcn & 4 & $0.376\pm0.042$ & $0.523\pm0.055$ & 0.332 & 0.574 \\
mhc\_lite\_gcn & 8 & $0.395\pm0.039$ & $0.544\pm0.050$ & 0.342 & 0.581 \\
\bottomrule
\end{tabular}
\end{table}

\textbf{Key findings:}
\begin{enumerate}
  \item \textbf{Val:} HC gains are learned ($+10\%$ Kendall, $\sigma\approx0.002$).
        Depth effect is neutral for GCN but significant for mHC (L4$\to$L8).
  \item \textbf{Test ID:} All models hit the same ceiling ($\approx82\text{--}83\%$
        Kendall). Hyper-connections are irrelevant on known distributions.
  \item \textbf{Test OOD:} Gains are substantial:
        HC-GCN $> $ mHC-GCN $\approx$ mHC-lite-GCN $>$ GCN.
        On top-K (NDCG@top5\%), HC and mHC nearly tie.
\end{enumerate}

%------------------------------------------------------------
\section{Stage 6 --- Residual Baseline (2026-03-10)}
%------------------------------------------------------------

\begin{table}[h]
\centering
\caption{Residual baseline --- OOD Kendall (5 seeds)}
\begin{tabular}{L{3.5cm} C{1.2cm} R{3.5cm} R{3cm}}
\toprule
Variant & $L$ & OOD Kendall & Relative gain \\
\midrule
gcn             & 4 & $0.305\pm0.046$ & --- \\
gcn\_residual   & 4 & $0.352\pm0.033$ & $+15.2\%$ \\
hc\_gcn         & 4 & $0.395\pm0.037$ & $+29.3\%$ \\
mhc\_gcn        & 4 & $0.380\pm0.049$ & $+24.4\%$ \\
\midrule
gcn             & 8 & $0.359\pm0.024$ & --- \\
gcn\_residual   & 8 & $0.369\pm0.035$ & $+2.6\%$ \\
hc\_gcn         & 8 & $\mathbf{0.411\pm0.040}$ & $+14.5\%$ \\
mhc\_gcn        & 8 & $0.389\pm0.046$ & $+8.4\%$ \\
\bottomrule
\end{tabular}
\end{table}

At L8: residual gives only $+2.6\%$; HC-GCN gives $+14.5\%$.
The hyper-connection gain is architectural, not reducible to
skip connections.

%------------------------------------------------------------
\section{Stage 7 --- Zero-Shot Real-Graph Transfer (2026-03-09/10)}
%------------------------------------------------------------

\subsection{Setup}
Models trained only on synthetic graphs evaluated zero-shot on:
\begin{itemize}[noitemsep]
  \item \textbf{Cora} (2708 nodes, 5278 edges, $\bar{d}\approx3.9$)
  \item \textbf{CiteSeer} (3327 nodes, 4552 edges, $\bar{d}\approx2.7$)
  \item \textbf{PubMed} (19717 nodes, 44324 edges, $\bar{d}\approx4.5$)
\end{itemize}

\subsection{Results --- Kendall $\tau$ (mean $\pm$ std, 5 seeds)}
\begin{table}[h]
\centering
\begin{tabular}{L{3.8cm} R{3cm} R{3cm} R{3cm}}
\toprule
Variant & Cora & CiteSeer & PubMed \\
\midrule
gcn\_l8           & $0.350\pm0.053$ & $0.067\pm0.118$ & $0.539\pm0.025$ \\
gcn\_residual\_l8 & $0.355\pm0.062$ & $0.069\pm0.080$ & $0.605\pm0.042$ \\
hc\_gcn\_l8       & $0.366\pm0.034$ & $0.056\pm0.029$ & $0.595\pm0.028$ \\
mhc\_gcn\_l8      & $\mathbf{0.405\pm0.036}$ & $\mathbf{0.151\pm0.083}$ & $0.598\pm0.030$ \\
mhc\_lite\_gcn\_l8 & $0.380\pm0.037$ & $0.125\pm0.085$ & $\mathbf{0.611\pm0.030}$ \\
\bottomrule
\end{tabular}
\end{table}

\subsection{Per-graph interpretation}

\textbf{Cora} ($\bar{d}\approx3.9$). All hyper-connected variants improve.
BC is not severely skewed. mHC-GCN leads on all metrics
($+15.5\%$ Kendall, $+43.8\%$ P@top5\%).
HC also gains ($+4.5\%$) without regression.

\textbf{CiteSeer} ($\bar{d}\approx2.7$ --- extreme sparsity).
Most informative graph. BC concentrates on a small number of bridge nodes.
HC-GCN \emph{regresses} $-16\%$ in Kendall vs.\ plain GCN,
while achieving best NDCG@10 ($+34\%$):
HC learns to route information toward bridges (correct top-K)
but destroys global ranking of other nodes.
The Sinkhorn constraint of mHC prevents this collapse:
$+127\%$ Kendall gain.

\textbf{PubMed} (19717 nodes).
Gains more moderate. GCN+residual is very competitive ($+12.2\%$).
mHC-lite leads ($+13.3\%$): full Sinkhorn overhead not worthwhile
at this scale; lighter variant offers better trade-off.
Precision@10 saturates at 0.5 for all models (non-discriminative).

\subsection{Routing collapse on sparse BC}

The HC regression on CiteSeer reveals a regime-dependent property:
\begin{center}
\begin{tabular}{L{5.5cm} L{5.5cm}}
\toprule
\textbf{HC (unconstrained)} & \textbf{mHC (Sinkhorn)} \\
\midrule
Best on synthetic BC (moderate skew) & Best on real citation graphs \\
High routing freedom & Proportional stream balancing \\
Collapses on sparse/skewed BC & No regression on any tested graph \\
\bottomrule
\end{tabular}
\end{center}

%------------------------------------------------------------
\section{Medium-Term Conclusions}
%------------------------------------------------------------

\subsection{What the project has established}

\begin{enumerate}
  \item \textbf{Hyper-connections improve inductive BC ranking.}
        Every HC-family variant outperforms GCN on OOD synthetic
        and on real-graph zero-shot transfer.
        The gain is architectural, not reducible to residual connections.

  \item \textbf{The gain is OOD-specific.}
        On in-distribution test graphs all models hit the same ceiling.
        Hyper-connections transfer knowledge; they do not add raw capacity
        on known distributions.

  \item \textbf{HC-GCN is the best variant on the synthetic protocol.}
        \[
          \text{HC-GCN}\ (0.406) > \text{mHC-GCN}\ (0.396)
          > \text{mHC-lite-GCN}\ (0.395) > \text{GCN}\ (0.360)
        \]

  \item \textbf{mHC-GCN is the most robust variant in zero-shot transfer.}
        It never regresses below the GCN baseline across any benchmark.
        HC-GCN collapses on very sparse, skewed-BC graphs (CiteSeer).
        The Sinkhorn constraint acts as routing regularization.

  \item \textbf{Deep scaling is not validated.}
        GCNII + mHC fails at L16. HC-family advantages concentrate at
        $L \leq 8$ for GCN backbones.
\end{enumerate}

\subsection{What cannot be claimed yet}

\begin{itemize}
  \item mHC does not universally outperform HC.
  \item Deep scaling ($L>8$) is not validated for any variant.
  \item The HC routing collapse is a single observation (CiteSeer);
        additional sparse graphs are needed to generalize it.
  \item No parameter-count-matched comparison exists;
        gains could partially reflect extra capacity.
\end{itemize}

\subsection{Recommended next steps}

\begin{enumerate}
  \item \textbf{BC spectrum analysis:} correlate BC distribution skewness
        (Gini, power-law exponent) of each graph with the HC vs.\ mHC
        relative performance, to test the routing-collapse hypothesis.
  \item \textbf{Internal matrix analysis:}
        inspect inter-stream transport matrices of HC-GCN-L8
        vs.\ mHC-GCN-L8 on CiteSeer to characterize the collapse
        geometrically.
  \item \textbf{Fair capacity comparison:}
        add a parameter-count table and capacity-matched baselines.
  \item \textbf{More sparse real graphs:}
        test on additional sparse social or collaboration networks
        to generalize the CiteSeer finding.
\end{enumerate}

\end{document}